\section{Lecture 10: inference 的总图}

本讲从一个端到端服务视角组织推理优化：先理解 prefill 与 generation 的算术强度和 KV cache，再区分有损 shortcut、带校验的无损 shortcut，以及处理动态请求的 continuous batching 与 PagedAttention。所有方法最终都要回到 latency、throughput、memory footprint 与质量四项账本。

\begin{figure}[H]
\centering
\includegraphics[width=0.92\textwidth]{inference-schema.png}
\caption{Lecture 10: inference 的整体结构：理解工作负载，有损 shortcut，无损校验，以及动态服务系统。}
\figsource{CS336 Spring 2026 Lecture 10 官方可执行讲义图像。}
\end{figure}

\begin{knowledgebox}{读图：inference-schema 如何串起整讲}
图里把本讲分成四块。第一块是 Understanding the inference workload：为什么推理和训练不同，为什么 generation memory-bound。第二块是 Taking shortcuts (lossy)：GQA/MLA/CLA、local attention、quantization、pruning/distillation 等可能改变模型或精度的做法。第三块是 Use shortcuts but double check (lossless)：speculative sampling 利用便宜 draft model 但保持目标模型精确分布。第四块是 Handling dynamic workloads：continuous batching 和 PagedAttention 解决真实请求到达、长度不同、共享前缀和 KV cache 碎片化。
\end{knowledgebox}

\begin{importantbox}{本讲的核心判断}
推理优化首先是\textbf{内存流量优化}，其次才是 FLOPs 优化。这里的 HBM 是 High Bandwidth Memory，即 GPU 上的高带宽显存；generation 阶段一次生成一个 token，不断从 HBM 读取模型权重和请求自己的 KV cache，常常 memory-bound。和 ZeRO/FSDP 训练状态分片不同，推理侧通常不维护 optimizer state；若要 shard 模型或 KV cache，目标是降低单卡 HBM 压力或提升吞吐，而不是同步训练梯度。
\end{importantbox}

\begin{knowledgebox}{讲义提醒：不要只报告单请求 latency}
推理系统的“快”必须绑定 workload。单请求、固定 batch、离线吞吐和在线到达流会得到不同最优策略；如果不报告 prompt/output 长度、并发度、硬件、精度与服务级目标，任何 tokens/s 数字都很难复现或比较。
\end{knowledgebox}

